Two research communities study the same failure in large language model
agents, and neither tests the other's work. Security researchers study the
context an attacker changes on purpose; machine learning researchers study
context that becomes wrong on its own. Either way the model cannot tell which
part of its context deserves trust. We measure that separation. From
\NCorpus{} papers read in full we recover \NPureTrackPairs{} confirmed
(defense, attack) pairs, of which only \NCrossCommunityPairs{}
(\PctCrossCommunity{}\%) cross between the communities, and one direction is
empty outright: \textbf{no security defense in our corpus has ever been
evaluated against a machine learning failure mechanism.} The dates show the
two sides reaching the same ground in turn rather than together. Of the eight
channels both have named a failure in, security names it first in five and
machine learning in none, by up to three years.

So we ran the missing experiment. We executed \NCampaignPairs{} (defense,
attack) pairs across \NCampaignImpl{} implementations and \NCampaignEpisodes{}
scored episodes, each from its authors' released code and weights and each
checked against an undefended model on the same items. Even where nine tenths
of the honest document survives, not one of them leaves the model better off
than no defense at all. The reason is a single mechanism. Each of these
defenses has one dial: how far to trust the retrieved document against how far
to trust the model's own memory. Turning it moves errors between the two
without creating correct answers. Trust the document more and the attacker,
who wrote part of that document, receives what memory would have supplied;
trust memory more and memory takes it back. The two movements almost exactly
cancel. Neither helps, because the attacker's text sits inside the very
document the defense is deciding whether to trust. One defense does resist. It
gives the retrieved text less weight when that text makes the model less
certain. So we wrote a passage that makes the model more certain. That turns
the weighting off, and cuts the defense's advantage over an undefended model
by 58\%.
